\section{Runtime refinement and trace adequacy}
\label{sec:refinement}

\subsection{Runtime-general refinement}

The refinement layer relates a concrete implementation machine to the
primitive event-graph machine.  It is deliberately runtime-general:
blockchain runtimes are intended backends, but the formal interface talks
about machines, states, observations, outcomes, and stochastic
transition kernels.

A stochastic refinement supplies a projection from implementation states
to specification states and proves that implementation transition laws
project to the specification transition laws.  Internal implementation
detail may be erased by the projection.  What matters is preservation of
the observable machine behavior and of the probability laws used by the
strategic analysis.

\subsection{Trace game surfaces}
\label{sec:ref:trace-surfaces}

Runtime adequacy is stated uniformly for a selected frontier game
surface.  A \(\mathsf{TraceGameSurface}\) packages:
\begin{itemize}\itemsep1pt
\item a strategy-indexed kernel game to preserve;
\item a primitive event-batch trace kernel realizing that game on the
  compiled specification machine;
\item a theorem that mapping primitive traces to utilities gives the same
  joint utility distribution as the game.
\end{itemize}
The canonical surfaces are behavioral frontier play, pure frontier play,
and product mixed-pure frontier play.

A \(\mathsf{TraceSpecEventBatchLaw}\) then gives a profile-indexed
specification event-batch law whose finite trace distribution is exactly
the trace kernel of the selected surface.  A
\(\mathsf{RuntimeTraceAdequacy}\) adds an implementation law family and
compatibility with the stochastic refinement.

\subsection{Transport theorem}
\label{sec:ref:transport}

The main transport statement is:
\begin{theorem}[Runtime trace adequacy]
\label{thm:runtime-trace}
Let \(g\) be a checked finite-domain program, let \(S\) be one of its
frontier trace game surfaces, and let an implementation machine refine
the primitive compiled machine.  If the implementation supplies a
strategy-indexed event-batch law compatible with a trace-adequate
specification law, then the implementation trace game has the same joint
utility distribution as \(S.\mathsf{game}\).
\end{theorem}

Under bounded-utility hypotheses, Nash equilibria of the selected
frontier surface pull back to Nash equilibria of the implementation
trace game.  The theorem is one-way.  A backend may expose additional
strategy space or implementation detail; refinement says that verified
specification equilibria remain valid under compatible implementation
laws, not that every implementation equilibrium is a specification
equilibrium.

\subsection{Identity and audited runtimes}

The identity runtime is the baseline: the primitive compiled machine
refines itself, and any trace-adequate specification law yields runtime
adequacy by reflexivity.  The audited runtime is a small generic example
in which the implementation executes an administrative audit tick before
specification event batches; the refinement erases those ticks.

These examples are smoke tests for the abstraction.  They do not bake any
particular backend into VegasCore.

\paragraph{Lean correspondence.}
Generic machine refinement is in \TT{Vegas.Machine.Refinement}.  The
program-facing trace-adequacy surface is in
\TT{Vegas.Core.Theorems.Refinement}: \TT{TraceGameSurface},
\TT{TraceSpecEventBatchLaw}, \TT{RuntimeTraceAdequacy},
\TT{behavioralFrontierTraceSurface}, \TT{pureFrontierTraceSurface},
\TT{mixedPureFrontierTraceSurface}, and the
\TT{implTraceGame\_udist\_surface} /
\TT{implTraceGame\_nash\_of\_surface\_nash} theorems.
